\emph{(1).} The signature of a continuous nondegenerate symmetric form on a connected base is constant: in a local frame its Gram matrix has continuous real eigenvalues, none of which vanishes, so the numbers of positive and negative ones are locally constant. This is the first step of the proof of Theorem \claimref{sheaf.conservation-degeneracy-absolute-form}, and it uses no tameness. Let $h$ be a strict completion standard outside $U$ and suppose it nondegenerate at every point of $M$. At $x, y \in M \setminus U$ its signature is $\operatorname{sig} G(x) + \operatorname{sig} K$ and $\operatorname{sig} G(y) + \operatorname{sig} K$, with the same $\operatorname{sig} K$ by strictness; these differ if $\operatorname{sig} G(x) \neq \operatorname{sig} G(y)$, against constancy on the connected $M$. So the set $D$ where $h$ degenerates is nonempty. It is closed, being the zero set of the determinant of a local Gram matrix. It lies in $U$, because on $M \setminus U$ the form is $G \perp K$ with both summands nondegenerate ($\Sig \subset U$). Applying the constancy to a path $\gamma$ from $x$ to $y$ (the pulled-back form on the pulled-back bundle over the interval) shows that $\gamma$ meets $D$. That theorem's tame case is the instance in which $U$ is a collar and the two values are $(p,q)$ and $(p-1,q+1)$.

\emph{(2).} A continuous $\beta$ with the stated properties exists by Urysohn's lemma, $\Sig$ and $M \setminus U$ being disjoint closed sets of a manifold; take $\beta = 0$ if $\Sig$ is empty. Work at one point, in a frame of $V$ orthonormal for $A$, and let $g$ be the symmetric matrix of $G$. The matrix of $h_{D}$ in the frame of $V \oplus V$ is
\[
\hat h = \begin{pmatrix} g & \beta \\ \beta & -g \end{pmatrix},
\qquad
\hat h^{2} = \begin{pmatrix} g^{2} + \beta^{2} & g\beta - \beta g \\ \beta g - g\beta & \beta^{2} + g^{2} \end{pmatrix}
= \begin{pmatrix} g^{2} + \beta^{2} & 0 \\ 0 & g^{2} + \beta^{2} \end{pmatrix},
\]
since $\beta$ is a scalar. The symmetric matrix $g^{2} + \beta^{2}$ is positive semidefinite, and positive definite when $g$ is invertible or $\beta \neq 0$; one of these holds at every point, because $g$ is invertible off $\Sig$ and $\beta > 0$ on $\Sig$. So $\hat h$ is invertible everywhere. With $J = \begin{psmallmatrix} 0 & 1 \\ -1 & 0\end{psmallmatrix}$ one has $J\hat h J^{-1} = -\hat h$, so the spectrum of $\hat h$ is symmetric about zero and its signature is $(n,n)$. Outside $U$, $\beta = 0$ and $h_{D} = G \perp (-G)$, with $K = -G$ on the second copy of $V$: nondegenerate there, of signature $(q,p)$ where $G$ has $(p,q)$, hence of one signature on $M \setminus U$ exactly when $\operatorname{sig} G$ is constant there. The form is the $h_{D}$ of clause (2) of Theorem \claimref{synthesis.wall-class-one-rung-obstruction}, where it is stated for tame walls with $\beta = 1$ near $\Sig$ and a signature of $K$ that changes; the computation above does not use tameness.

\emph{(3).} (1) and (2). The obstruction is that the locally constant function $\operatorname{sig} G$ on $M \setminus U$ is not constant.

\emph{(4).} In (a), $x^{2} > 0$ for $x \neq 0$: signature $(1,0)$ on both sides, corank one at $x = 0$, and $x^{2}$ vanishes there to second order, so the section $\delta$ of (T3) vanishes on $\Sig$. In (b) the eigenvalues are $x$ and $-x$: signature $(1,1)$ for $x \neq 0$ and corank two at $0$. In (c), $x^{2} + y^{2} > 0$ off the origin, and $\det\begin{psmallmatrix} x & y \\ y & -x\end{psmallmatrix} = -(x^{2} + y^{2})$ with trace zero, so the signature is $(1,1)$ off the origin; $\Sig$ is a point, not a hypersurface, and $M \setminus \Sig$ is connected. In each case $\operatorname{sig} G$ is constant on $M \setminus \Sig$, hence on $M \setminus U$ for every $U$, and (2) applies. For (a) the smaller form has determinant $-x^{2} - \beta^{2} < 0$ where $x \neq 0$ or $\beta \neq 0$, so signature $(1,1)$ everywhere, and equals $x^{2} \perp (-1)$ where $\beta = 0$.

\emph{(5).} For $x \neq 0$ the form $x\,\id$ is definite, positive for $x > 0$ and negative for $x < 0$; the two signatures differ, so (1) applies with any $U \ni 0$ whose complement meets both half-lines. Their residues are $4$ and $-4 \equiv 4 \pmod 8$, and the two Clifford algebras are isomorphic by the classification (\claimref{std-algebra.periodicity}). A change of $p - q$ modulo $8$ is in particular a change of $(p,q)$. In the tame corank-one case one eigenvalue changes sign, $(p,q) \to (p-1,q+1)$, and $p - q$ changes by $2$, the clock mismatch of the stored theorems.

\emph{(6).} At such a point $V_{x}$ with the form $G_{x}$ is a quadratic space with radical of dimension $r$ inside the nondegenerate space $(E_{x}, h_{x})$, which is therefore a nondegenerate completion of it; by Fact \claimref{odd-extensions.imported-minimal-completion-degenerate-algebra-rule} its dimension is at least $(n - r) + 2r = n + r$, and since the radical is a totally isotropic subspace of dimension $r$ of a nondegenerate space, that space has at least $r$ positive and $r$ negative directions.

\emph{Reading.} The conjecture's strict reading is the statement of (1) with ``one side'' and ``the other side'' for two exterior regions; (1), (4) and (5) give the three assertions of the Reading, and the last sentence is clause (2) of Theorem \claimref{synthesis.wall-class-one-rung-obstruction} with (6).
